\chapter{反射与序列化}
\label{ch:20}

当多个工具都需要字段名称、类型或声明上的附加信息时，可以通过元数据避免重复维护映射。注解提供声明信息，反射按支持的接口查询或访问对象，序列化再把值编码为外部表示。

本章用 Route 和 Profile 区分这些职责，再把字段访问放入真实输入验证流程。减少重复代码之后，权限、范围和对外字段选择仍由应用明确决定，不能全部交给一个反射标记。

\section{显式暴露元数据}

先给一个函数附上路由信息。path 表示希望关联的路径，level 表示应用可能解释的级别；这些数据和函数的输入、输出并不是同一件事。函数即使被标记，也仍按普通规则执行自己的函数体。

@annotation 定义这些附加信息的形状，使用注解时再提供具体值。下面保留最小的直接调用，先确认声明信息与执行动作的区别；真正的路由分派还需要应用去消费元数据。

\par
\medskip
\begin{lstlisting}[style=chali]
@annotation Route(path: String, level: int32 = 2)

@Route("/health", level: 1)
int64 health()
{
    return 200;
}

void main()
{
    println("{}", health());
}
\end{lstlisting}
\par
\medskip

\par
\medskip
\begin{lstlisting}[style=out]
200
\end{lstlisting}
\par
\medskip

\subsection*{代码详解}

@annotation 定义 Route 的元数据形状：path 必须提供 String，level 是 int32 且缺省为 2。health 上的 @Route 提供路径并显式把 level 覆盖为 1，这些值属于声明信息。

main 直接调用 health，函数体 return 200，再由 println 输出，所以结果来自普通函数执行。代码没有创建 HTTP 服务，也没有读取 Route 来分派请求；200 在这里就是整数，不能仅凭它与 HTTP 状态码同名便认定已经发送了网络响应。

\section{按字段读取与写入}

反射访问仍然有具体类型。字段名不存在、读取类型不匹配或不允许写入时，应处理接口的失败结果。下面通过字段元数据读取、修改业务值；调用位置不能把反射当成绕过只读或所有权规则的入口。

\par
\medskip
\begin{lstlisting}[style=chali]
import std::reflect;

@Reflect
struct Profile
{
    String name;
    int64 score;
};

void main()
{
    Profile profile;
    profile.name = "原名";
    profile.score = 80;
    assert(std::reflect::setField<Profile, String>(&profile, "name", "新名").unwrap());
    assert(std::reflect::getField<Profile, String>(&profile, "score").isErr());
    assert(std::reflect::getField<Profile, int64>(&profile, "missing").isErr());
    int64 score = std::reflect::getField<Profile, int64>(&profile, "score").unwrap();
    println("{} {}", profile.name, score);
}
\end{lstlisting}
\par
\medskip

\par
\medskip
\begin{lstlisting}[style=out]
新名 80
\end{lstlisting}
\par
\medskip

\subsection*{代码详解}

@Reflect 为 Profile 暴露所需模型信息。setField 的两个类型实参分别指定模型和写入值类型，\texttt{\&profile} 提供目标对象，字段名 name 定位字符串字段。unwrap 与 assert 在这个已知成功的例程中确认操作结果，name 随后成为“新名”。

后两次查询故意失败：score 不能按 String 读取，missing 也不是真实字段，因此用 isErr 验证。最后用正确的 int64 查询 score，取得 80，输出“新名 80”。反射写入只修改指定字段，没有把失败的查询变成默认值；真实外部输入应显式处理 Result，避免对未知结果直接 unwrap。

\section{序列化边界}

JSON 等外部数据需要先解析，再检查结构与字段约束。类型正确的整数仍可能是不合理的年龄，存在的班级编号仍可能超出当前教师权限。对外响应选择允许显示的字段，密码摘要等内部字段不应因为参与反射就被自动导出。

\section{元数据与生成的职责}

把可重复的类型映射交给反射，把业务决策保留在清楚的函数边界。第 25 章用数据库与页面说明这些边界：字段校验、关联记录和操作权限需要分别验证。

\subsection*{把字段信息放进完整的数据流程}

以上 Profile 示例包含三个独立条件：类型暴露了反射信息，字段名指向实际字段，请求的值类型与字段类型相容。score 存在并不意味着可以把它作为 String 读取；missing 不存在也不意味着 score 的元数据出了问题。处理失败时保留这些区别，才能判断是接口调用写错、外部字段名错误，还是模型没有提供所需信息。

在真实应用中，可以把处理分成读取外部表示、验证数据、修改业务对象、选择输出字段几个步骤。比如修改成绩时，先检查输入能否解析成整数，再检查分数范围、学生与课程是否存在，以及当前账号是否有权修改，最后才写入对象和持久化。反射减少字段映射代码，但不会替应用决定检查顺序或事务边界。

注解同样需要解释者。Route 示例只定义并使用声明上的数据，直接调用 health 的输出不能证明 HTTP 路由已经注册。要提供路由服务，还需要程序读取或使用相应信息、建立分派，并处理请求。区分“声明了元数据”和“执行了业务动作”，可以避免把注解误当成会自动启动的功能。

运行期反射与编译期生成也应分别理解。当前 getField、setField 示例让程序在运行时按字段名请求具体类型的访问；它不展示编译器遍历异构字段并为每个字段生成静态代码。后者涉及尚未交付的宏与编译期反射扩展。已有派生指令只按其现行约定使用，不能从存在反射 API 推断任意自定义派生已可用。详细接口见\href{https://github.com/rabbitc999/ChaLi/blob/master/doc/19-reflection.md}{注解与反射}。

\subsection*{小结}

从声明元数据到业务动作，中间仍需要解释、查询、验证和处理失败。先确认类型是否暴露信息，再确认字段及请求类型，最后执行允许的读写；输出外部数据时另选公开字段与编码协议。

运行期反射与编译期派生有各自支持范围，不因名称相近而互相包含。下一章跨入宿主库，类型、布局和所有权约定将更需要明确记录。

\section{习题}

\begin{enumerate}

\item 注解是否会自己执行路由？

\item 读取错误类型的字段如何处理？

\item 为何不能直接序列化账号完整结构体？

\item 列出反射不能代替的两项业务检查。

\end{enumerate}

\subsection*{参考解答与验收要点}

\begin{enumerate}

\item 不会，需要程序解释并分派。

\item 检查返回结果，不把失败当成正常值。

\item 内部敏感字段可能泄露；明确选择响应字段。

\item 例如班级存在性和当前教师的数据访问范围。

\end{enumerate}

下一章：\hyperref[ch:21]{第 21 章}。

\section{番外：榜上有名，方可查验}

改写自《封神演义》姜子牙封神的相关情节。

子牙登台，展开封神榜，众魂各候敕命。闻太师拍案道：“贫道虽死，字段尚全，先给我查个类型！”守榜官翻了半日，道：“尊驾没有登记可反射的模型信息，单凭生前威名，读不出字段偏移。”

申公豹在旁笑曰：“道友留步！我有画像一张，像得十分，何必元数据？”守榜官问画像之中字段叫什么、各是什么类型。申公豹袖中只有一张脸，连 \texttt{name} 都没声明。

又有人榜上有名，取出整数，却执意按字符串使用，口称“反正最后都要显示”。子牙将打神鞭一横：“元数据告诉你是什么，不负责把你想要的东西变出来。”

守榜官欲将一切失败都填成“不在榜上”，子牙不许：未登记、字段缺失、类型不合，是三桩事。反射能力本在通常构建中，莫因一次查询为空，就宣称整个天庭裁掉了查榜台。

封到最后，众人各有职司。申公豹仍问：“既有注解，业务规则可否全免？”子牙合榜道：“榜上写你司掌雷部，不代表你能把用户的年龄填成一万岁。”

当日天庭最大的收获，是一份模型供众处查询；最大的节约，是终于不用让三十六位雷将各自手抄同一张字段表。

\subsection*{技术对照}

注解提供元数据，不自动实现业务逻辑。模型未暴露、字段缺失或类型不匹配，与编译器裁剪掉反射能力是不同情况；高级入口的 panic 边界按实际构建约定处理。
